\clearpage
\begin{summarybox}
\textbf{一句话记住：}
\textbf{把难算的函数换成好算的式子；若前面的项会抵消，就继续留后面的项，直到不会影响答案为止。}

\textbf{五条公式速记（$x$ 接近 $0$，三角函数用弧度）：}
\[
\begin{aligned}
 e^x&\approx1+x+\frac{x^2}{2},\\
 \ln(1+x)&\approx x-\frac{x^2}{2}+\frac{x^3}{3},\\
 \sin x&\approx x-\frac{x^3}{6}+\frac{x^5}{120},\\
 \cos x&\approx1-\frac{x^2}{2}+\frac{x^4}{24},\\
 \tan x&\approx x+\frac{x^3}{3}.
\end{aligned}
\]

\textbf{三步使用法：}
\begin{enumerate}[leftmargin=*]
\item \textbf{先检查：} $x$ 是否接近 $0$，角度是不是弧度，原函数有没有定义。
\item \textbf{再看抵消：} 分子减掉了哪些项，就把近似式多写到下一项。
\item \textbf{最后查误差：} 求极限时，省略部分除以分母必须趋于 $0$；
估小数时，误差要小到不改变最终保留的位数。
\end{enumerate}

\textbf{两个“看到就想起”的结果：}
\[
 \lim_{x\to0}\frac{1-\cos x}{x^2}=\frac12,
 \qquad
 \lim_{x\to0}\frac{\tan x-\sin x}{x^3}=\frac12.
\]
它们可以只用单位圆与三角恒等式说明，特别适合先理解，再记忆。

\textbf{额外提醒：}
$\ln(1+x)\le x-x^2/2+x^3/3$ 对一切 $x>-1$ 都成立；
这比“两个式子差不多”更有用，因为还告诉了我们\textbf{大小方向}。

\medskip
\textbf{新增：一张图记住二项式主线。}
\[
 \boxed{\begin{gathered}
 (1+x)^n\quad(n=0,1,2,\ldots)\\
 \downarrow\ \text{改看倒数，先用乘法核对}\\
 \frac1{1+x}=1-x+x^2-x^3+\cdots\\
 \downarrow\ \text{比较求导后的各项}\\
 \ln(1+x)\approx x-\frac{x^2}{2}+\frac{x^3}{3}
 \end{gathered}}
\]
\[
 \boxed{\left(1+\frac{x}{m}\right)^m\xrightarrow[m\to\infty]{}
 e^x\approx1+x+\frac{x^2}{2!}+\frac{x^3}{3!}+\cdots}
\]
再往后，把 $x$ 换成 $ix$ 并按两部分分组，
可以\textbf{联想}正弦、余弦的系数规律；
但这一步必须额外知道 $e^{ix}=\cos x+i\sin x$，
\textbf{不当作高中课内的独立证明}。
正切则由 $\tan x=\sin x/\cos x$ 得到。

\textbf{三个绝不能忘记的条件：}
$n=-1$ 不再有固定的最后一项；
$1/(1+x)$ 的无限项写法先要求 $|x|<1$；
$\ln x$ 的这个近似要看 $x$\textbf{接近 $1$}。
\end{summarybox}
